\partsection{F.5}{Contributions, AI use, and known limitations}\label{F:5}

\subsubsection*{Contribution table}

From the merged pull requests and closed issues of \path{beambeambeam/claimflow} on GitHub, collected by \path{scripts/collect_team.py} on 15 November 2026, after the release.
An equal number of commits is not required; the two students split the work as follows.
Dependabot pull requests (\#10, \#22) are not counted.

% Generated by scripts/gen_team_table.py from team.json. Do not edit.
\begin{center}\small
\begin{tabular}{@{}p{4.8cm}p{5.2cm}p{5.2cm}@{}}
\toprule & \textbf{67070501005 Kiatisakk} & \textbf{1045 beambeambeam} \\ \midrule
Merged pull requests authored & 7 & 17 \\
Closed issues assigned & 6 & 11 \\
PRs reviewed (teammate's) & 8 & 7 \\
\quad with an approval & 8 & 6 \\
\quad with changes requested & 3 & 0 \\
Pull requests, numbers & \#18, \#20, \#21, \#59, \#60, \#68, \#69 & \#8, \#9, \#11, \#12, \#13, \#14, \#15, \#17, \#19, \#23, \#35, \#36, \#56, \#57, \#58, \#63, \#65 \\
Issues, numbers & \#1, \#2, \#3, \#27, \#28, \#29 & \#5, \#6, \#7, \#24, \#25, \#26, \#30, \#31, \#32, \#33, \#34 \\
\bottomrule\end{tabular}\end{center}

{\footnotesize\begin{longtable}{@{}p{0.9cm}p{7.4cm}p{2.6cm}p{3.9cm}@{}}
\toprule \textbf{PR} & \textbf{Title} & \textbf{Author} & \textbf{Review by the other student} \\ \midrule \endhead
\#8 & feat: set up bun monorepo with web and api apps & beambeambeam & none recorded \\
\#9 & chore: add repo health files & beambeambeam & none recorded \\
\#11 & build: proxy web api requests to the api in dev & beambeambeam & none recorded \\
\#12 & chore: add mattpocock agent skills & beambeambeam & none recorded \\
\#13 & feat: validate api env config at startup & beambeambeam & none recorded \\
\#14 & docs: add converted project specification documents & beambeambeam & none recorded \\
\#15 & build: replace t3 env with varlock in web & beambeambeam & none recorded \\
\#17 & ci: speed up api job & beambeambeam & none recorded \\
\#18 & docs: add software requirements specification & Kiatisakk & commented \\
\#19 & refactor: rename app to claimflow & beambeambeam & none recorded \\
\#20 & docs: add feature inventory and traceability matrix & Kiatisakk & approved, commented \\
\#21 & docs: add system-level software design specification & Kiatisakk & approved \\
\#23 & docs: add sprint plan, feature specs, tickets, and sprint 1 contracts & beambeambeam & approved, changes requested \\
\#35 & docs: rename sprint branches and add review rules & beambeambeam & approved \\
\#36 & feat: add walking skeleton with contract, database, and health checks & beambeambeam & approved, dismissed \\
\#56 & feat: log in and show the current user & beambeambeam & approved, changes requested \\
\#57 & docs: move issue tracking to GitHub Issues & beambeambeam & approved \\
\#58 & feat: log out, session expiry, and bans & beambeambeam & approved \\
\#59 & feat: add route permissions and the route inventory test & Kiatisakk & approved, commented \\
\#60 & feat: submit a claim with its rules and field errors & Kiatisakk & approved, commented \\
\#63 & feat: list my claims and open a claim with its policy & beambeambeam & approved, commented \\
\#65 & feat: review queue, approve and reject a claim, and the Sprint 1 journey & beambeambeam & approved, changes requested \\
\#68 & chore: prepare the Sprint 1 release & Kiatisakk & approved, commented \\
\#69 & chore: release Sprint 1 & Kiatisakk & approved \\
\bottomrule\end{longtable}}


\begin{center}\small
\begin{tabular}{@{}>{\bfseries}p{2.7cm}p{13.6cm}@{}}
\toprule
Kiatisakk (1005) & Requirements and design documents: SRS (\#18, issue \#1), feature inventory and traceability matrix (\#20, issue \#2), system-level SDS (\#21, issue \#3). Implementation: route permissions and the route inventory test, F-01/04 (\#59, issue \#27); submit a claim, F-02/01 and F-02/02, with the claim rules, the schema and the New Claim form (\#60, issues \#28, \#29). Reviews of Beam's PRs \#23, \#35, \#36, \#56, \#57, \#58, \#63, \#65, with changes requested on \#56 and \#65. The release: the check of \texttt{sprint-1} against the Product DoD (\#68) and the release PR (\#69). This report (issue \#4). \\
beambeambeam (1045) & Repository, CI and tooling setup (\#8, \#9, \#11 to \#15, \#17, \#19, issue \#7), the sprint plan and the Sprint 1 contracts (\#23, issues \#5, \#6), the sprint workflow rules (\#35), moving issue tracking to GitHub Issues (\#57). Implementation: walking skeleton with contract, database and health checks, F-01/01 (\#36, issue \#24); log in and the current user, F-01/02 (\#56, issue \#25); log out, session expiry and bans, F-01/03 (\#58, issue \#26); My Claims and the claim detail, F-03/01 and F-03/02 (\#63, issues \#30, \#31); the review queue, approve and reject, and the Sprint 1 journey with Playwright, F-04/01 to 03 (\#65, issues \#32, \#33, \#34). Reviews of Kiatisak's PRs \#18, \#20, \#21, \#59, \#60, \#68, \#69, and the merge of the release. \\
\bottomrule
\end{tabular}
\end{center}

Pair work: none declared in any Sprint 1 pull request (the ``Pair work'' section of PR \#60 says ``None'').
\todo{1045: confirm the split above and the roles of the second student}

\subsubsection*{AI-use declaration}

\textbf{Student 67070501005 (Kiatisak Markmeeshap).} Claude Code (Anthropic) was used throughout Sprint 1 with Claude Fable 5.2 and Claude Opus 5.5 as orchestrator, and Claude Sonnet 5.5 subagents as implementers and reviewers.
Tasks: implementation of tickets under test-driven development (the failing test is committed first, see F.3), drafting of code-review comments on the teammate's pull requests, and tooling for this report (the capture, cut and check scripts under \path{docs/reports/sprint-1/scripts/}).
Every result is gated by the repository checks (formatting, lint, tests, coverage) and reviewed by a person before it is committed, pushed or posted; the review comments in F.4 were read and posted by the student.

\textbf{Student 1045 (Supawit Marayat).} \todo{1045: AI tools used and the tasks they were used for}

\subsubsection*{Known limitations, failed and skipped tests}

Written from the contracts, the pull request descriptions and the open issues at the release commit \texttt{ac5bce5}.
Issues \#61, \#66, \#67 and \#70 track the open work; the two small known gaps of PR \#60 have no issue of their own.

\begin{itemize}\setlength\itemsep{2pt}
\item \textbf{The E2E specs do not run in CI} (issue \#66). F01-TEST-24, F02-TEST-21, F03-TEST-12, F04-TEST-23 and F04-TEST-24 pass on a local seeded stack (F.3), but CI has no Playwright job. For that reason the ``pass in CI'' box of \S7 in the F-01, F-02 and F-03 contracts is left open with a known-gap line (Part E); F-04's Playwright item allows a recorded run.
\item \textbf{The SDS and the code differ in three places} (issue \#67): two \texttt{List} methods where \S3.1 has one, history actions stored as strings, and empty audit fields for the session. SDS revision 0.12 records them as as-built notes; \#67 also holds the smaller review follow-ups of PR \#65. The review of the release PR \#69 added \#70: none of its items blocks the release or changes an acceptance criterion; the one that touches behaviour is the F-04 \S2.2 403 page, which shows only when the next read of the claim also returns 403.
\item \textbf{Receipt files are not part of Sprint 1} (issue \#61). A line item has a \texttt{Receipt reference} text field; uploads come later.
\item \textbf{Layout shift in error rows.} In a New Claim line item that shows errors, the columns shift right (PR \#60, known gaps).
\item \textbf{The F-02 E2E test covers the happy path only;} the rule errors are covered by the API and component tests (PR \#60).
\item \textbf{Manual checks} were run on \texttt{29fd6fb}, the last code commit before the release, whose code is identical to \texttt{ac5bce5} (the two later commits change \texttt{README.md} and \path{docs/} only): F01-TEST-18 (no tracked \texttt{.env}; gitleaks reports 11 hits, all Go module checksums; \texttt{.env.example} matches \texttt{.env.schema}), F01-TEST-19 (required CI checks on PR \#65), and F04-TEST-19 (restart check: \texttt{req.anan}'s claims, items and history byte-identical before and after), which was run again on \texttt{ac5bce5} (F.3). Results are in \path{docs/traceability.md}.
\item \textbf{Failed and skipped tests on the release commit.} Go: 0 failed, 0 skipped of 123 top-level tests and 152 subtests. Web: 130 of 130 pass in CI on \texttt{ac5bce5}; on the author's machine a full local run times out on 19 tests under load, and rerun file by file only the first test of \texttt{auth-routing.test.tsx} still fails there (F.3). E2E: 5 of 5 pass locally. No test was disabled, skipped or weakened to obtain these results.
\end{itemize}
